
\chapter{Управление отказами}\label{ch:decline-management}

Покупатель на~странице оплаты видит одно и~то же сообщение: платёж
не~прошёл. За~этим сообщением стоят разные коды ответа: один
допускает повтор через паузу, другой говорит о~проблеме маршрута
при~исправном счёте клиента, третий запрещает повторять операцию
с~теми же реквизитами. \highlight{Различить эти случаи нужно до~того,
как система начнёт повторять заведомо безнадёжные попытки}.

\needspace{4\baselineskip}
\section{Структура отказа}\label{sec:decline-anatomy}

Дихотомия \enquote{мягкий~/ жёсткий отказ} остаётся эвристикой
продавца. Visa делит коды отказов (decline response codes) на~категории
\enquote{Issuer will never approve}, \enquote{Issuer cannot approve at this
time}, \enquote{Data quality} и~\enquote{Generic}; Mastercard
отдельно запрещает дополнительные авторизационные запросы по~CNP для части
кодов отказов~\cite{visa-core-rules,mc-tpr}. Продавцу
важно одно: какое действие разумно после конкретного кода.

Мягкий отказ (\textit{soft decline}) означает, что транзакция отклонена по
вр\'{е}менной или устранимой причине: временный лимит, поведенческий
фильтр, недоступность сервиса. К~этой группе чаще
всего относят \texttt{51} Insufficient Funds, \texttt{61}~Exceeds
Withdrawal Amount Limit, \texttt{65} Exceeds Withdrawal Frequency
Limit и~часть общих (\textit{generic}) кодов вроде \texttt{05}, если
схема вообще допускает повтор и~ТСП остаётся в~пределах лимитов
повторов~\cite{visa-core-rules}.

Жёсткий отказ (hard decline) означает, что ту же операцию с~теми же
данными повторять бессмысленно или прямо запрещено правилами.
Надёжные примеры~--- \texttt{04} Pick Up,
\texttt{14} Invalid Account Number\footnote{В~актуальной редакции
Visa Core Rules код~\texttt{14} отнесён к~категории~1 (Issuer will
never approve); операционно он эквивалентен жёсткому отказу~--- повтор
с~теми же реквизитами бессмысленен, нужен новый PAN.}, \texttt{41}
Lost Card и~\texttt{43} Stolen Card. Они либо указывают на~неверные
реквизиты, либо на~карту, с~которой продолжать операцию нельзя,
и~поэтому требуют другого платёжного реквизита или ручной
проверки~\cite{visa-core-rules,mc-tpr}.

\texttt{54} Expired Card
для Visa относится к~категории Data quality, поскольку схема прямо
трактует его как случай, где нужно перепроверить платёжную информацию.
У~Mastercard правило сформулировано иначе: для CNP ТСП не должен
инициировать дополнительные авторизационные запросы по~той же операции
с~тем же PAN и~сроком действия~--- повторять тот же запрос без~новых
реквизитов нельзя~\cite{visa-core-rules,mc-tpr}.

Технические и~маршрутные коды стоят отдельно. \texttt{91} Issuer
Unavailable и~\texttt{96} System Malfunction говорят о~доступности
эмитента или маршрута, не~о~состоянии счёта. Для них схема допускает
повторные попытки в~пределах лимитов; сначала нужно убедиться, что
проблема в~доступности пути, сама карта в~этих
кодах ни~при~чём. \texttt{01} Refer to Card Issuer слепо повторять
нельзя: обычно это знак, что нужен иной канал или дополнительное
действие со~стороны держателя либо эмитента~\cite{visa-core-rules}.

\begin{table}[H]
  \begingroup
  \scriptsize
  \setlength{\tabcolsep}{4pt}
  \renewcommand{\arraystretch}{1.2}
  \begin{minipage}{\textwidth}
    \begin{tabularx}{\textwidth}{@{}Y Y Y@{}}
      \toprule
      \cellcolor{Bad!16}\textbf{Не повторять ту же операцию}
      &
      \cellcolor{Warn!18}\textbf{Нужны новые данные / пауза}
      &
      \cellcolor{AccentA!14}\textbf{Технический / другой путь}
      \\
      \midrule
      \cellcolor{Bad!4}\textbf{04 Pick Up}\newline стоп; запросить другой реквизит
      \newline\textbf{14 Invalid Account}\newline проверить PAN
      \newline\textbf{41 Lost Card}\newline карта утеряна; стоп
      \newline\textbf{43 Stolen Card}\newline карта украдена; стоп
      &
      \cellcolor{Warn!5}\textbf{54 Expired Card}\newline обновить срок
      действия~/ PAN; в~Mastercard CNP не повторять тот же PAN+expiry
      \newline\textbf{51 Insufficient Funds}\newline пауза; без цикла повторов
      \newline\textbf{61 Exceeds Limit}\newline проверить сумму и~лимит
      \newline\textbf{62 Restricted Card}\newline смотреть контекст и~правила схемы
      \newline\textbf{65 Freq. Limit}\newline дождаться сброса лимита
      \newline\textbf{05 Do Not Honor}\newline общий отказ; только в~пределах лимитов схемы &
      \cellcolor{AccentA!5}\textbf{91 Issuer Unavailable}\newline
      проблема доступности маршрута или эмитента
      \newline\textbf{96 System Malfunction}\newline технический сбой;
      сначала зафиксировать финальный статус
      \newline\textbf{01 Refer}\newline нужен иной канал или действие банка
      \newline\textbf{таймаут}\newline подтвердить финальный статус,
      затем контролируемый повтор \\
      \bottomrule
    \end{tabularx}
    \par\smallskip
    \tablenote{Точная схема сложнее этой таблицы: Visa делит
      коды отказов на~Category~1/2/3/4,
      а~Mastercard отдельно запрещает часть CNP retries. Здесь
      показан только первый операционный вопрос:
      прекращать операцию, обновлять данные или работать с~доступностью
      маршрута.}
  \end{minipage}%
  \endgroup
  \caption{Рабочая классификация кодов ответа по~следующему действию.}\label{tab:decline-codes}
\end{table}

\section{Значение кодов ответа}\label{sec:decline-response-codes}

Код ответа эмитента не объясняет причину отказа: за~одним и~тем же
кодом стоят десятки причин, и~без дополнительных данных истинную
причину не узнать.

\texttt{05} Do Not Honor формально означает \enquote{не проводите
операцию} и~почти ничего не объясняет. За~ним скрываются антифрод,
дневной лимит, несоответствие поведенческому профилю клиента,
временная блокировка в~приложении или внутренняя политика эмитента
по~конкретному MCC. Часть причин временная и~исчезает сама, часть
постоянная, поэтому \texttt{05} трактуют с~осторожностью. У~Visa
общие коды отказов (generic decline response codes) сведены в~отдельную
Category~4 и~должны использоваться только там, где не подходит более
точное значение, поэтому \texttt{05} читается как отказ
без расшифровки~\cite{visa-core-rules}. Эмитент намеренно не раскрывает
причину через конкретный код: детальный код~--- информация, полезная
атакующему. Если \texttt{51} точно означает нехватку средств,
злоумышленник с~украденными картами отбросит \enquote{пустые}
и~сосредоточится на~\enquote{богатых}. Возвращая \texttt{05}, эмитент
лишает атаку перебором краденых учётных данных
(\textit{credential stuffing}) информации о~состоянии счёта.

\texttt{51} Insufficient Funds тоже не сводится к~буквальному прочтению.
Для дебетовых карт он часто означает нехватку средств на~счёте,
для кредитных~--- превышение лимита, а~у~некоторых эмитентов~--- универсальную
форму мягкого отказа, когда раскрывать точную причину они не хотят.

\texttt{54} Expired Card в~таблице категорий кодов отказа Visa Core
Rules относится к~классу Data quality, а~в~Visa Token Service (VTS)
PAN Lifecycle прямо предусмотрены обновление PAN, срока действия
и~обновления Visa Account Updater (VAU). В~потоках с~повторными
списаниями и~сохранёнными реквизитами карт после \texttt{54}
проверяют не возможность повтора, а~факт получения обновлённых
реквизитов. У~Mastercard ту же операционную задачу закрывает
Automatic Billing Updater (ABU), который публично позиционируется
как способ держать платёжные реквизиты актуальными и~избегать
отказов~\cite{visa-core-rules,visa-vts,visa-vau,mc-abu,mc-bill-pay-solutions}.

\texttt{91} и~\texttt{96}~--- технические коды о~доступности эмитента
или маршрута, не~о~состоянии карты или счёта. В~Visa оба кода попадают
в~классы, где повторы разрешены в~пределах лимитов, заданных схемой.
Если такие коды доминируют в~статистике, проблема в~сети,
маршрутизации или доступности
хоста~\cite{visa-core-rules}.

Мягкие отказы, связанные с~аутентификацией, образуют отдельную
категорию. Код~\texttt{65} в~исходной ISO~8583/Visa-семантике
означает Exceeds Withdrawal Frequency Limit, когда нужно дождаться
сброса лимита частоты; в~европейском контексте Strong Customer
Authentication (SCA, усиленная аутентификация плательщика)
Mastercard тот же \texttt{65} используется в~качестве soft decline,
по~которому эмитент готов одобрить операцию после повторной
аутентификации по~протоколу 3-D~Secure (3DS)\,---\,см.~гл.~\ref{ch:3ds}.
Поток после такого отказа устроен так. Шлюз получает \texttt{65},
распознаёт его как допускающий повтор с~3DS-аутентификацией
(\textit{retryable-with-3DS}\,---\,внутренний класс шлюза, а~не тег схемы),
инициирует аутентификацию: сообщение Authentication Request
(AReq) уходит от~3DS-сервера ТСП на~Directory Server (DS)
платёжной схемы и~далее на~Access Control Server (ACS) эмитента.
По~результатам аутентификации шлюз
получает Electronic Commerce Indicator (ECI, индикатор уровня
аутентификации) и~Cardholder Authentication Verification Value
(CAVV, криптограмма подтверждения держателя) и~повторяет
авторизационный запрос с~данными аутентификации. Эмитент видит
тот же PAN, ту же сумму, но~с~подтверждением SCA, и~одобряет. Слепой
повтор того же запроса даст лишь цикл отказов; повтор с~3DS
возвращает шанс на~одобрение~\cite{mc-tpr}.

\practicenote{%
  Хороший мониторинг отказов сегментирует коды ответа не~по~двум
  группам, а~хотя бы по~трём-четырём~--- жёсткие, мягкие финансовые,
  мягкие поведенческие и~технические. Сводная панель с~одним числом
  доли отказов скрывает структуру. 15\,\% отказов, половина из~которых
  приходится на~технические коды, означают совсем другую задачу, чем
  те~же 15\,\% при~преобладании необратимых отказов.

}
\subsection*{Офлайновые коды POS-терминала: Z1, Z3, Y1, Y3}\label{sec:decline-offline-rc}

Не все коды ответа приходят от~эмитента. Если терминал по~правилам
EMV или по~настройкам конфигурации сам принимает решение о~судьбе
транзакции~--- одобрить или отклонить, не~отправляя её в~сеть,~---
он генерирует код собственного формата. Эти коды не~входят в~ISO~8583,
видны только в~чеке POS и~в~локальном журнале; эмитент о~такой
операции вообще не~узнаёт~\cite{emvco-book4}.

Решение терминала задаётся набором Issuer Action Code (IAC)~--- это
записанные на~карте битовые маски реакций EMV-приложения на~результаты
Terminal Risk Management. IAC-Denial указывает, по~каким сигналам
приложение карты требует отклонить операцию офлайн, IAC-Online~---
по~каким сигналам обязательно уйти на~авторизацию к~эмитенту,
а~IAC-Default~--- какое решение принять, если уйти онлайн не
удалось~\cite{emvco-book3}.

\begin{itemize}
  \item \texttt{Z1}~--- \textit{offline declined}: терминал отклонил операцию
    самостоятельно, не~пытаясь связаться с~эмитентом. Типично~---
    срабатывание условий IAC-Denial, провал офлайн-аутентификации
    данных, превышение прикладного лимита.
  \item \texttt{Z3}~--- \textit{no online, declined}: терминал хотел уйти
    онлайн (требование IAC-Online), но связи не~было, и~по~правилам
    IAC-Default операцию пришлось отклонить.
  \item \texttt{Y1}~--- \textit{offline approved}: терминал одобрил операцию
    локально, не~отправляя её на~авторизацию. Возможно только
    если конфигурация терминала разрешает офлайн-одобрение
    и~сумма не~превысила floor limit.
  \item \texttt{Y3}~--- \textit{no online, approved}: попытка уйти онлайн
    провалилась, но IAC-Default позволил одобрить операцию офлайн.
\end{itemize}

В~российской рознице \texttt{Z1} и~\texttt{Y3} встречаются редко:
floor limit у~большинства терминалов равен нулю, и~любая операция
уходит на~эмитента. \texttt{Z3} и~\texttt{Y3}, напротив, всплывают
при~разборе инцидентов с~падением канала связи и~помогают локализовать
проблему между терминалом, эквайером и~сетью.

\section{Повторы после отказа: время, частота, каналы}\label{sec:decline-retry}

На~повторах расходятся интересы продавца и~платёжной схемы: продавцу
нужна конверсия, схеме~--- бережное расходование авторизационной
ёмкости. При~миллионах подписочных ТСП даже два лишних запроса в~месяц
на~клиента создают значимую нагрузку, и~схема ставит лимиты повторов
там, где их польза близка к~нулю.

Правила схем задают внешние условия повторов:
какие коды вообще можно повторять, а~какие нельзя. Visa Core Rules
сводят коды отказов в~таблицу категорий с~разными лимитами
повторов для ТСП; Mastercard для CNP отдельно запрещает дополнительные
авторизационные запросы по~той же транзакции после ряда кодов отказа,
включая \texttt{04}, \texttt{14}, \texttt{41}, \texttt{43} и~\texttt{54}~\cite{visa-core-rules,mc-tpr}.

Схемы почти не~говорят ТСП
\enquote{повтори через 17~минут}; они говорят \enquote{не~делай слепой
цикл повторов и~не~выходи за~лимиты}. Если отказ связан с~нехваткой
средств или лимитом, серия одинаковых запросов не~создаёт деньги и~не
расширяет лимит. Если код технический, бессмысленная очередь повторов
только зашумляет логи и~повышает риск попадания под мониторинг схемы.

\highlight{Перед каждым повтором продавец отвечает на~два вопроса:
допускает ли схема повтор этого кода или повтор требует новых данных
(срок действия, токен, 3DS), а~если повтор допустим~--- какой интервал
и~какой потолок безопасны для конкретного продавца, PSP и~категории
ТСП.} Mastercard на~публичных продуктовых страницах продаёт ту~же
логику как \textit{intelligent retry strategies}~\cite{mc-bill-pay-solutions}.

Отдельный фактор~--- состояние холда эмитента (\textit{auth-hold})
после первой авторизации (см.~раздел~\ref{sec:tx-authorization}). Если первая
попытка вернула мягкий отказ уже после того, как эмитент успел
поставить hold (например, тайм-аут на~стороне эквайера), повтор
по~тому же реквизиту может прийти эмитенту как дубль и~быть
отклонён, а~суммарная блокировка средств у~держателя~--- удвоиться
до~истечения периода действия hold по~MCC. Грамотная логика повторов
проверяет состояние первой попытки и~при~необходимости снимает блокировку
сообщением \textit{reversal} перед выпуском второй авторизации.

По~допустимости повтора коды ответа удобно разложить на~классы:

\begin{itemize}
  \item \textbf{Жёсткий отказ} (\texttt{04}, \texttt{14},
    \texttt{41}, \texttt{43})~--- карта недействительна, изъята
    или украдена. Повтор запрещён, нужен новый платёжный
    инструмент.
  \item \textbf{Мягкий финансовый} (\texttt{51})~--- нехватка
    средств. Повтор допустим, но~не~ранее чем через 24--48~часов,
    когда баланс может измениться; цикл с~интервалом в~минуты
    бесполезен.
  \item \textbf{Мягкий аутентификационный} (\texttt{65},
    \texttt{1A})~--- требуется SCA. Повтор идёт через 3DS;
    повтор того же запроса бесполезен.
  \item \textbf{Технический} (\texttt{91}, \texttt{96})~---
    эмитент или маршрут недоступен. Повтор допустим через
    15--60~минут; каскад через другого эквайера эффективен
    только если альтернативный маршрут ведёт к~другому хосту
    эмитента.
  \item \textbf{Непрозрачный} (\texttt{05})~--- причина неизвестна.
    Серия идентичных запросов повышает риск попадания
    под~мониторинг схемы.
\end{itemize}

Visa и~Mastercard описывают токенизацию и~оптимизацию одобрений
как способы дать эмитенту более чистые данные и~убрать лишние
шаги авторизации, поэтому повтор через сетевой токен или через
более удачный маршрут не~равен простому повтору запроса
по~PAN~\cite{visa-vts,visa-tokenization-security,mc-acquirers-psps,mc-pop}.

Логика повторов сокращает бессмысленные попытки; долю отказов
снижают на~уровне реквизитов, маршрутов и~аутентификации, и~эти
две задачи решаются разными инструментами.

\section{Сетевой токен и~PAN: дополнительный контекст доверия}\label{sec:decline-tokens}

Если продавец работает через сетевой токен (без~хранения PAN),
эмитент видит реквизит с~привязанными ограничениями. Visa Token Service
описывает токен как отдельный цифровой идентификатор, с~которым
эмитент задаёт token variables, определяет авторизованных
token requestors и~запрашивает token details, включая device и~risk
information. Токен может быть
ограничен конкретным мобильным устройством, ТСП категории
\textit{e-commerce} или количеством покупок~\cite{visa-vts,emvco-payment-tokenisation-guide}.

В~публичных материалах Visa называет рост доли одобрений (approval
rate) для токенизированных транзакций порядка нескольких процентных
пунктов; Mastercard на~страницах для acquirers и~PSP пишет о~3--6~п.\,п.\
роста доли одобрений у~tokenized digital
payments~\cite{visa-tokenization-security,mc-acquirers-psps}. Цифры
принадлежат самим схемам и~каждому продавцу не~гарантируются, но общий
вывод устойчив: токен даёт эмитенту более чистую картину, чем обычный
PAN.

Эффект токенизации виден в~жизненном цикле реквизитов.
Visa PAN Lifecycle обновляет PAN, срок действия и~выполняет VAU updates,
а~Mastercard Automatic Billing Updater подаётся как способ держать
платёжные реквизиты актуальными и~избегать отказов. В~потоках
с~сохранёнными реквизитами карт \texttt{54} Expired Card чаще означает
устаревший реквизит у~продавца, а~не~у~эмитента~\cite{visa-vts,visa-vau,mc-bill-pay-solutions}.

Миграция на~токены откладывается по~инерции: если текущий поток
работает, переход откладывают до~момента, когда доля отказов и~оттока
становится заметной. Многие продавцы не~управляют токенами напрямую
и~зависят от~темпа своего PSP.

Подробнее об~архитектуре сетевых токенов см.~раздел~\ref{sec:token-anatomy}.

\section{Маршрутизация через эквайера}\label{sec:decline-routing}

Платежи доходят до~эмитента по~нескольким маршрутам, и~выбор пути
меняет исход авторизации. В~правилах Mastercard есть отдельные разделы
про \textit{authorization routing} и~\textit{domestic transaction
routing}, а~в~публичных материалах для acceptance community
\textit{routing} и~\textit{payment optimization} вынесены
в~самостоятельную линейку
продуктов~\cite{mc-tpr,mc-acquirers-psps,mc-pop}. Выбор одного PSP
описывает лишь точку входа, не~весь путь до~эмитента.

У~разных эквайеров различаются локальное покрытие, подключения, качество
формирования сообщения и~сервисы оптимизации. Сетевая оптимизация
одобрений работает на~этих различиях: подбирает удачный маршрут
авторизации или дополняет сообщение недостающими
данными~\cite{mc-pop}.

Страна BIN и~страна ТСП определяют, транзакция внутренняя (\textit{domestic})
или международная (\textit{cross-border}); внутренний маршрут через локального
эквайера даёт эмитенту знакомый контекст и~часто лучше конвертирует.
MCC влияет на~доступные программы \textit{interchange} и~на~набор
правил антифрода эмитента. Валюта транзакции и~валюта расчёта
важны, если платформа поддерживает DCC. Историческая доля одобрений
на~конкретном эквайер--BIN--MCC сочетании служит входом для
выбора маршрута в~реальном времени.

Принцип выбора маршрута~--- функция ожидаемого дохода
$P(\text{approve} \mid \text{route}) \times \text{amount} -
\text{cost}(\text{route})$, где маршрут с~более высокой комиссией
оправдан только при~достаточно высоком $P(\text{approve})$. В~зрелых
системах роль такой функции играет модель ML, переобучаемая на
исторических авторизациях, а~в~простых~---
скоринговая (\textit{score}) таблица по
BIN~$\times$~MCC~$\times$~страна.

Различие \textit{cascade routing} (каскадный перебор эквайеров) и
\textit{intelligent routing} (выбор маршрута по~предсказанной вероятности
одобрения) как двух шаблонов маршрутизации, корректная обработка
задержавшегося ответа первого эквайера и~ограничения каскада по~типам
отказов разобраны в~разделе~\ref{sec:gw-routing}. Операционное следствие
для этой главы: если все эквайеры в~каскаде вернули отказ,
транзакция завершается с~последним полученным кодом, повторно прогонять
её через тех же эквайеров запрещено правилами повторов схем,
а~правильное действие~--- показать ошибку или предложить другой
платёжный метод.

Правила схем, локальной маршрутизации и~interchange запрещают
маскировать географию ТСП или искусственно подбирать маршрут под
выгодную тарифную модель. Легитимная цель маршрутизации~---
повысить долю одобрений и~сократить операционные
потери~\cite{visa-core-rules,mc-rules,mc-tpr}.

\section{Доля одобрений по~категориям ТСП}\label{sec:decline-benchmarks}

Одно и~то же значение доли одобрений~--- например, 85\,\%~--- интерпретируется
по-разному в~зависимости от~категории ТСП, географии, типа карты
и~модели оплаты.

Сетевые правила и~мониторинги считают approval performance не~одним
универсальным числом: Mastercard отслеживает долю одобрений в~CNP-
и~международных контекстах через Merchant Advice Code (MAC),
Authorization Optimizer и~связанные мониторинговые программы.
Сравнивать все потоки одной агрегированной метрикой нельзя~\cite{mc-tpr}.

Диапазоны доли одобрений устаревают быстрее, чем выходит книга,
потому что они зависят от~категории ТСП, страны, доли 3DS,
качества токенизации, текущих волн мошенничества и~даже
календарного периода. Универсальные ориентиры \enquote{нормальной}
доли одобрений по~категориям ТСП бесполезны: опубликованные числа
не~учитывают разрезов конкретного продавца. Конкретные
проценты для своего потока сверяются по~данным PSP или
эквайера. Системные просадки относительно собственной
исторической нормы по~категории ТСП указывают на~конкретные сбои:
неправильный повтор, отсутствие 3DS на~мягкий отказ, невыверенный
маршрут по~BIN или устаревшие реквизиты
в~\textit{card-on-file}.

Сравнивать розницу с~сервисами путешествий, подписки с~разовыми
интернет-платежами, а~внутренний поток с~международным в~одной колонке
бессмысленно. В~подписочных операциях (\textit{recurring}) добавляется контекст
сохранённого реквизита. Если цепочка согласия не~оформлена по
правилам SCF, эмитент видит обычный запрос CNP; продолжение
уже данного согласия в~таком запросе не~считывается~\cite{visa-scf,mc-tpr}.

\practicenote{%
  Долю одобрений измеряйте раздельно по~сегментам: новые и~вернувшиеся
  клиенты, карты с~3DS и~без неё, PAN и~токен, внутренний рынок и
  международные операции.
}
